\documentclass[10pt,a4paper]{amsart}
\usepackage[margin=19mm]{geometry}
\usepackage{amsmath,amssymb,mathtools}
\usepackage[scr=rsfs,cal=euler]{mathalfa}
\usepackage{xfrac}
\usepackage[unicode,hidelinks,bookmarks=false]{hyperref}
\newcommand{\ie}{i.e.\ }
\newcommand{\cf}{cf.\ }
\newcommand{\resp}{respectively\ }
\newcommand{\Subsection}{\subsection}
\providecommand{\nobreakdash}{\nobreak\hbox{-}}
\setlength{\emergencystretch}{2em}
\multlinegap0pt
\usepackage{bm}
\usepackage{bbm}
\usepackage{dsfont}
\usepackage{xspace}
\usepackage{cancel}
\usepackage{mathtools}
\usepackage{amsthm}
\makeatletter
\@ifundefined{theorem}{\newtheorem{theorem}{Theorem}}{}
\makeatother
\title[Prime Quadruplets and Jump Conditions on Arithmetic Function]{Prime Quadruplets and Jump Conditions on Arithmetic Functions}
\author{Himaghna Roy Choudhury, Shicheng Wei}
\date{Source: arXiv:2606.10331v4 (CC BY 4.0)}
\begin{document}
\maketitle

\noindent\emph{Source and licence.} Prime Quadruplets and Jump Conditions on Arithmetic Functions. Version arXiv:2606.10331v4 (CC BY 4.0), distributed under the Creative Commons Attribution 4.0 International licence, as recorded in the arXiv licence metadata for this exact version and shown on the abstract page https://arxiv.org/abs/2606.10331v4. Full rights notice: licenses/arxiv-2606.10331-CC-BY-4.0.txt.\par\smallskip

\noindent\textit{Editorial scope.} Counted unit: the Theorem whose source label is \texttt{thm:semiprime-char} in the article (source0.tex lines 102--107, source label \texttt{thm:semiprime-char}), delivered together with the article's own complete original proof of it (source0.tex lines 109--182, one \texttt{proof} environment of 74 lines). No mathematics is written, repaired or re-derived: statement and proof are the article's own text line for line. Required context retained uncounted: none. Every definition, display and label of those lines is retained unchanged. Omitted from this package and not redistributed: 1--101, 108--108, 183--545. Nothing inside a retained excerpt is omitted or altered: every retained body is delivered exactly as the article's lines are written, so no omission receipt of any kind accompanies this package. Citations: the retained excerpts carry no citation command at all, so no bibliography entry of the article is transcribed and no reference is redistributed. Reference adaptation: none was needed and none was made; every reference inside the delivered unit resolves inside it, so no unresolved reference or citation is produced. Printed slips kept literal: no printed word, symbol, hypothesis or number of the counted statement or of its proof was altered. Layout and class: the article's own class is \texttt{amsart} with packages including geometry, mathtools, booktabs, float, longtable, xcolor, hyperref, url, amsmath, amsthm, amsfonts, amssymb; the wrapper is the lane's pinned \texttt{amsart} editorial wrapper at 10pt on A4, which restates the theorem-like environments the retained text needs and adds the article's own macro definitions that the retained text needs (none), with extra wrapper packages (bm, bbm, dsfont, xspace, cancel, mathtools). The delivered numbering is the unit's own. Assets: no retained span contains a figure environment, a graphics-inclusion command, a PDF-page inclusion, a table or a TikZ picture; no image file is copied into this package, the package declares no asset dependency and no image file is redistributed with it. Licence: version arXiv:2606.10331v4 of this article is distributed under the Creative Commons Attribution 4.0 International licence, as recorded in the arXiv licence metadata for this exact version and shown on the abstract page https://arxiv.org/abs/2606.10331v4. A full rights notice is delivered at licenses/arxiv-2606.10331-CC-BY-4.0.txt. Characters: every retained excerpt is pure ASCII, so the delivered engine, XeLaTeX with its default Latin Modern text font, reports no missing character.
\begin{theorem}\label{thm:semiprime-char}
Suppose $n = pq$ and $n+12 = p'q'$ are both products of two distinct primes (with $p < q$ and $p' < q'$). Then both jump conditions
\[
\varphi(n+12) - \varphi(n) = 12, \qquad \sigma(n+12) - \sigma(n) = 12
\]
hold if and only if $p \ge 5$, $q = p+8$, and $\{p', q'\} = \{p+2,\, p+6\}$. That is, $(p,\, p',\, q',\, q) = (p,\, p+2,\, p+6,\, p+8)$ is a prime quadruplet.\end{theorem}

\begin{proof}
($\Leftarrow$) Using $\varphi(rs) = (r-1)(s-1)$ and $\sigma(rs) = (r+1)(s+1)$ for distinct primes $r$ and $s$, we get
\begin{align*}
\varphi(n) &= (p-1)(p+7) = p^2 + 6p - 7, \\
\varphi(n+12) &= (p+1)(p+5) = p^2 + 6p + 5, \\
\sigma(n) &= (p+1)(p+9) = p^2 + 10p + 9, \\
\sigma(n+12) &= (p+3)(p+7) = p^2 + 10p + 21.
\end{align*}
By subtracting, we get $\varphi(n+12) - \varphi(n) = 12$ and $\sigma(n+12) - \sigma(n) = 12$ from just polynomial identities in terms of $p$.

($\Rightarrow$) For any squarefree semiprime $rs$ with $r$ and $s$ distinct primes,
\[
\varphi(rs) = (r-1)(s-1)=rs - (r+s) + 1, \qquad \sigma(rs) = (r+1)(s+1) = rs + (r+s) + 1.
\]
Applying this to $n = pq$ and $n+12 = p'q'$:
\begin{align*}
\varphi(n)      &= pq - (p+q) + 1                       &&= n - (p+q) + 1, \\
\varphi(n+12)   &= p'q' - (p'+q') + 1                   &&= (n+12) - (p'+q') + 1, \\
\sigma(n)       &= pq + (p+q) + 1                       &&= n + (p+q) + 1, \\
\sigma(n+12)    &= p'q' + (p'+q') + 1                   &&= (n+12) + (p'+q') + 1.
\end{align*}
Subtracting the first from the second and the third from the fourth:
\begin{align*}
\varphi(n+12) - \varphi(n) &= \bigl[(n+12) - (p'+q') + 1\bigr] - \bigl[n - (p+q) + 1\bigr] = 12 + (p+q) - (p'+q'), \\
\sigma(n+12)  - \sigma(n)  &= \bigl[(n+12) + (p'+q') + 1\bigr] - \bigl[n + (p+q) + 1\bigr] = 12 + (p'+q') - (p+q).
\end{align*}
Both expressions equal $12$ if and only if
\begin{equation}\label{eq:sum-eq}
p + q = p' + q'.
\end{equation}
Now, let $S = p+q = p'+q'$. Then, $p, q$ are the roots of $X^2 - SX + n = 0$, and $p', q'$ are the roots of $X^2 - SX + (n+12) = 0$. The discriminants are
\[
u^2 := S^2 - 4n = (q - p)^2 \ge 0, \qquad v^2 := S^2 - 4(n+12) = (q' - p')^2 \ge 0,
\]
so
\begin{equation}\label{eq:disc-diff}
u^2 - v^2 = 48.
\end{equation}

We now seek integer pairs $(u, v)$ with $u > v \ge 0$ satisfying \eqref{eq:disc-diff}. We first write $48 = (u-v)(u+v)$, which leads us to notice that $u-v$ and $u+v$ must have the same parity, since their sum $2u$ is even. 

If both factors are odd, their product is odd, which contradicts the fact that $48$ is even. If both are even, then let $u - v = 2a$ and $u + v = 2b$ with $a < b$ and $ab = 12$. The factorizations $(a, b) \in \{(1, 12),\, (2, 6),\, (3, 4)\}$ yield
\[
(u, v) = (a + b,\, b - a) \in \{(13, 11),\, (8, 4),\, (7, 1)\}.
\]

We can verify them directly: $7^2 - 1^2 = 48$, $8^2 - 4^2 = 48$, $13^2 - 11^2 = 48$. These are the only integer solutions with $u > v \ge 0$.

Now recall $S = p + q$. We consider separate cases for whether $p = 2$:

\emph{Case 1: $p, q$ both odd.} Then $S = p+q$ is even, and $u = q - p$ must satisfy $S - u = 2p$ (even), forcing $u$ even. Of the three solutions, only $(u, v) = (8, 4)$ has an even $u$. Hence
\[
q - p = 8, \quad q' - p' = 4.
\]
\begin{samepage}
Combined with $p' + q' = p + q = 2p + 8$, this gives us a linear system in $p', q'$:
\[
p' + q' = 2p + 8, \qquad q' - p' = 4.
\]
Adding and subtracting yields $2q' = 2p + 12$ and $2p' = 2p + 4$ respectively. Hence
\[
p' = p + 2, \qquad q' = p + 6.
\]
\end{samepage}
Since $p, \, p', \, q', \, q$ are prime by hypothesis, we know that $(p, p+2, p+6, p+8)$ is a prime quadruplet. In addition, $p \ge 5$, because the smallest prime quadruplet starts at $p = 5$ (if $p = 3$, then it requires $3, 5, 9, 11$ all to be prime, but $9 = 3^2$ is not).

\emph{Case 2: $p = 2$.} Then $n = 2q$ is even, so $n + 12 = p'q'$ is also even. Since $p', q'$ are both prime and $2$ is the only even prime, either $p'$ or $q'$ must be $2$. By the convention $p' < q'$, we have $p' = 2$, which leads to
\[
q' = \frac{n+12}2 = \frac{2q+12}2 = q + 6.
\] 
Condition~\eqref{eq:sum-eq} gives $2 + q = 2 + q'$, i.e., $q = q' = q + 6$, which is impossible. (Alternatively, $\varphi(2q) = q - 1$ and $\varphi(2(q+6)) = q + 5$, with difference $6 \neq 12$.) So this case yields no solutions.
 
This exhausts the cases.
\end{proof}

\end{document}
